\section{Introduction}
\label{sec:intro}

Agents built on language models spend much of their budget talking to each other. The tokens exchanged grow with every hand-off, and much of that traffic is redundant \cite{zhang2024cut}. The cost appears as latency and energy and, on edge hardware, where memory is fixed and the longest prompt sets the batch size, as a limit on how many requests can be served at all \cite{zheng2025areview,shkolnikov2026agentmemory}. Communication efficiency is now a design problem in its own right \cite{aei2026survey}.

The literature offers three levers. The first is to \emph{say less}: prune redundant messages from the agent graph \cite{zhang2024cut}, send only the semantic packets whose expected contribution exceeds their cost \cite{dong2026bandmas}, or compress the prompt \cite{mu2023gist,ge2024icae,cheng2024xrag,jiang2023llmlingua}. The second is to \emph{reuse computation}: persist and share caches so that a receiver does not re-read what it has already read \cite{shkolnikov2026agentmemory}. The third is to \emph{send a latent message} instead of text, handing the receiver hidden states, a key--value cache or an embedding with no decoding into words. The training-free LatentMAS reports 70.8 to 83.7\% fewer tokens and 4 to 4.3 times faster inference than text on nine benchmarks \cite{zou2025latent}; cache-to-cache reports a 2.5-fold reduction in latency and 3.1 to 5.4\% higher accuracy than text \cite{fu2025cachetocache}. Other designs follow the same pattern \cite{pham2024ciphers,du2026enabling,zheng2025thought,rossi2026latentcacheflow,flamant2026bicameral,jin2026agentprimitives}, and a recent framework organises them by what is transmitted and how the receiver absorbs it \cite{liu2026beyondtokens}. Systems have begun to choose among these levers adaptively, for instance between token and cache transmission according to bandwidth \cite{dai2026tokenkv}.

All of this work assumes that the agents are language models, so that the sender's hidden state or cache is something another transformer can be trained to absorb. The agents that sense the world usually are not transformers. A wearable or bedside classifier is a convolutional or recurrent network trained for classification, and nothing about its internal vector is shaped for a language model. The one heterogeneous exception we know of connects different families of vision--language models through a shared visual channel \cite{liu2026visionwormhole}, but its senders are still language models. For a sender of this kind it is open whether a latent channel is feasible, whether it is worth building, and what it costs in inspectability.

Neighbouring work does not settle the question. Projecting a signal encoder into a language model's embedding space is established \cite{jin2024timellm,babu2025modality,zhao2025ecgchat,zhang2025sensorlm}, but those systems are built to make a language model understand a signal and are judged on accuracy. They do not treat the projection as a channel whose cost is compared with a text channel on the same task. Comparisons of latent and text communication also tend to use the text baseline that is easiest to beat. Listing every event verbatim is a poor choice when the sender can filter, which is the first lever above. Without a filtered baseline, a saving in tokens says little about whether a latent channel is worth building.

This paper compares the levers for a non-transformer sender. A small learned adapter maps the 32-dimensional context vector of a convolutional--recurrent heartbeat classifier, together with three scaled side inputs, to four virtual tokens per event in the input-embedding space of a frozen, 4-bit Gemma~4 E4B model. We compare it with three text interfaces, among them a sender-side filter, each also measured with prefix caching. The task is to report the most urgent triage tier among 1 to 50 consecutive heartbeats. Electrocardiography is the case study: deep networks already match cardiologists on it \cite{hannun2019cardiologist} and language models encode clinical knowledge \cite{singhal2023large}, so pairing a sensing agent with a reasoning agent is a natural design. The task tests \emph{communication fidelity}. The correct answer is a fixed rule applied to the sender's outputs, so a wrong answer is a failure of the channel and not of clinical knowledge. A three-line rule also solves the task, and we do not claim otherwise: the language model adds nothing to the decision here. It stands in for the general-purpose receiver of an agent system, a model shared by many senders and tasks for which no one writes a rule per task, and the benchmark isolates the channel into it. The experiment therefore measures what each interface costs and loses, under a protocol fixed before the results existed, and the results tables report the rule itself as the ceiling. We ask three questions.
\begin{itemize}
\item \textbf{RQ1 (cost).} Does the adapter reduce the tokens, time and memory of the hand-off relative to text interfaces, and under which workloads?
\item \textbf{RQ2 (accuracy).} Does it preserve the accuracy of the downstream decision?
\item \textbf{RQ3 (recoverability).} How much of each event's content can be recovered from the virtual tokens, so that the cost of opacity is measured and not asserted?
\end{itemize}

The answers do not all favour the latent channel. It keeps prompt processing flat as events are added, while text grows by about 60 tokens per event, and it is more accurate than calibrated compact text at moderate window sizes. A prompt listing only the abnormal events, the filter, is more accurate than the adapter and as fast for a single request. Its length depends on how many abnormal events a window contains, whereas the adapter's depends only on the number of events, and this decides the comparison once requests are batched. We report the result in this form because a system designer needs a rule for which lever to pull.

The contributions are as follows.
\begin{enumerate}
\item A communication-fidelity benchmark for the interface choice of a non-transformer sender, comparing the three efficiency levers (filter, reuse, latent) with a text baseline that includes a sender-side filter, on a task whose answer is fixed so that every error is the channel's.
\item A latent channel from such a sender into a frozen language model, extended from one event per call to fifty, in which the language model is untouched and only 0.88 million adapter parameters are trained.
\item A controlled evaluation of cost, accuracy and recoverability against a pre-specified analysis plan, including serving under batching and prefix caching, replication on a second sender architecture, and an external confirmatory test that confirmed both the adapter's advantage over calibrated text and the filter's advantage over the adapter (Section~\ref{sec:incart}).
\item A recoverability protocol that decodes each event's label, tier, heart rate, interval and abnormal-run status from its slot of virtual tokens, separating what the channel carries from what the receiver uses.
\item A rule for choosing between a filtered text interface and a latent one, with the analysis plan, its change log and code to reproduce every table.
\end{enumerate}

Section~\ref{sec:related} places the work among the three levers and the literature on heterogeneous senders. Section~\ref{sec:system} describes the sender, the adapter and the decoding scheme, and Section~\ref{sec:protocol} the evaluation protocol. Sections~\ref{sec:results} and~\ref{sec:discussion} report and discuss the results, with limitations.
